\section{Limitations}

\begin{enumerate}\itemsep2pt

\item \textbf{Depth and density are not separable.} Every analysis in this paper is
affected by the fact that sequencing depth proxies tissue density. Where the choice is
consequential we report both adjusted and unadjusted estimates rather than selecting one.

\item \textbf{No per-cell malignant label.} We do not assign malignant status to any cell,
spot or domain. The copy-number score is reported as a continuous quantity throughout, and
no downstream statistic multiplies it into any other measurement.

\item \textbf{Domain identity is method-sensitive.} BANKSY, RCTD composition and
marker-based clustering resolve 7, 6 and 5 domains respectively, with an adjusted Rand
index of 0.225--0.247 against an RCTD-aligned partition. The seven domains are a working
description, not a unique decomposition.

\item \textbf{The deconvolution reference bounds what deconvolution can report.} The
39-subtype reference has no exhaustion programme, so exhausted T cells are reported under a
naive label. T-cell composition in this paper is a composition in the reference's
vocabulary; functional states not present in the reference are invisible to it, and were
measured separately by signature projection.

\item \textbf{The $K^{*}=7$ partition failed its pre-specified stability gate.} Across
seedings the cross-seed adjusted Rand index never reached 0.90. We report the partition as
exploratory and do not use it as a primary result; the consequences are confined to
descriptions of domain composition. The companion spatial-coherence criterion we also
applied is passed by every configuration tested, including one using no spatial
information, and therefore carries no weight in this cohort.

\item \textbf{One clinical stage is represented by one section.} The Normal arm of the
spatial cohort consists of a single Visium section; MIA is represented by four. Estimates
for these stages should not be quoted on their own. Nuclei are more balanced (92,337
normal nuclei across 23 patients), so the imbalance is specific to the spatial arm.

\item \textbf{Two domain identities were revised during the study.} The D3 domain was
first described as an immune infiltrate, then as a hypoxic/invasive programme, and finally
as ECM/interstitial-enriched after depth matching showed the first two readings to be
driven by depth. Earlier names for this domain should not be used.

\item \textbf{The marker panel for one analysis was hand-curated.} A 22-gene ECM panel
used in the early spatial analyses overlaps the depth-balanced marker set by only 14
genes; the remaining eight were added by us. A sensitivity analysis shows the finding is
unchanged when those eight are removed, or when any single gene is removed, but the panel
itself is not drawn from a single published source.

\item \textbf{Stage comparisons may be contaminated by sample type.} Patient-internal
pairing removes between-patient variation but not the systematic difference between a
precursor lesion and an invasive carcinoma, which is present in every paired comparison.
The consistency of a signal across 22 patients does not protect against this, since every
patient contributes the same kind of contrast.

\item \textbf{The prognostic analysis is exploratory.} Domain marker signatures were used
without a signed pre-registration of the parameterisation, and one domain's result changes
substantially when its signature is replaced by a depth-balanced version. Only the
direction and its robustness are reported for that domain.

\item \textbf{The compound candidates are untested.} They act in cancer cell lines, have
not been assessed in lung fibroblasts or primary tissue, and have no experimental
validation. They are candidate perturbations.

\item \textbf{The gene-level candidate list rests on one entry.} After the specificity gate,
the direction filter and the compartment decomposition, only two perturbations rank in the
top 100 across twenty or more of the 31 subtype axes, and one of the two fails the
direction check. The surviving gene-level claim therefore rests on a single perturbation
(\emph{STAT1} knockdown), which we report as an observation rather than as a target.

\item \textbf{The ligand layer produced nothing.} NicheNet passed only five ligands, none of
which yielded a ligand--target link, so no ligand-level conclusion can be drawn from this
cohort. The failure is traceable to a thin sender panel and is not evidence that no ligand
effect exists.

\item \textbf{The perturbation libraries are not independent.} The gene-level and
compound-level libraries share the query signature, and the three compound scoring schemes
share both the query and the library. Agreement between them is a robustness check, not
independent corroboration. The two libraries are reported side by side and are not used to
validate one another.

\item \textbf{The pooled gene-perturbation query produced no signal.} Its positive control
passes, so the null reflects the query rather than the method, but we cannot distinguish
between ``this cohort's pooled disease axis is not represented in the library'' and ``no
single-gene perturbation of this magnitude exists in the library''. Splitting the query by
compartment recovers a ranking, so the null is specific to the pooled parameterisation
rather than to the library.

\item \textbf{The developmental analysis is ambiguous by construction.} The score we use
is derived from the number of detected genes, so a high value in invasive disease is
equally consistent with greater transcriptional disorder. Two readings fit the same
numbers and we cannot separate them with the available labels.

\item \textbf{The imaging analysis is bounded by resolution, not by biology.} The
H\&E-based classifier ceiling is set by the model architecture against the available image
scale, and the model does not separate the precursor stages from one another at all. The
negative result concerns what this pipeline can extract at this resolution, not what is
present in the tissue; a finer colour source would be needed to test the latter.

\item \textbf{No orthogonal validation.} No result in this paper has been reproduced in an
independent cohort, with a different perturbation library, or experimentally.

\end{enumerate}
